\chapter{Introduction}
% =============================================================================

\section{Motivation}

Cross-modal retrieval - the task of matching relevant items in one modality given a query in another - has become an increasingly important task in modern machine learning, underpinning systems for image search, audio captioning, text-to-video matching, and an expanding catalogue of multimodal applications. The dominant model couples two encoders through a contrastive objective trained on large corpora of paired examples: CLIP \cite{radford2021clip} learns from approximately 400 million image--text pairs, ALIGN \cite{jia2021align} from roughly a billion noisy image--text pairs, and CLAP \cite{elizalde2023clap} from on the order of $10^6$ audio--text pairs. 

Although contrastive methods have shown incredible results in matching heterogeneous modalities, they are limited by (i) fundamentally data hungry approaches that require large, paired datasets \cite{Yuan2021Multimodal} and (ii) an inability to distinguish between highly similar samples due to broad alignment \cite{rösch2024enhancingconceptualunderstandingmultimodal}.

For instance, despite the perceptual link between sound and image, there is no widely adopted contrastive foundation model that spans the two modalities natively. AudioCLIP~\cite{guzhov2021audioclip} does extend CLIP to a third audio branch via a tri-modal contrastive loss, but it remains a relatively small-scale experiment and has not supplanted the workarounds that practitioners more commonly reach for,such as distilling CLIP into an audio encoder~\cite{wu2021wav2clip}, or bridging CLIP and CLAP through their shared text modality~\cite{wang2023cmcr,wang2023exmcr}. 
The state of the art in audio--visual generation~\cite{cheng2025mmaudio} and in self-supervised audio--visual representation learning~\cite{alwassel2019xdc,morgado2021avid} sidesteps this problem by jointly modeling audio and video through attention mechanisms.

This thesis investigates an alternative. Following the recent Platonic Representation Hypothesis \cite{huh2024platonic} we ask to what extent we can align audio and visual representations in both completely unsupervised and low supervision settings, leveraging the internal geometric structure of pretrained encoders in each modality. Concretely, given a set of items in modality $A$ embedded by one model and a set of items in modality $B$ embedded by an independently trained model, we study whether the relative geometry of points within each space - how images relate to other images compared with how texts relate to other texts - is sufficient to recover meaningful correspondences across the two spaces.

\section{Platonic Representation Hypothesis}

The premise that intra-modal geometry alone may suffice for cross-modal alignment is supported by a recent line of work on representation convergence. The Platonic Representation Hypothesis \cite{huh2024platonic} conjectures that as neural networks scale, the representations they learn converge towards a shared statistical model of reality, irrespective of training data, modality, or objective. Ziyin et al. \cite{ziyin2025proof} provide a constructive proof of a ``perfect'' Platonic regime for embedded deep linear networks under SGD, and López-Cardona et al. \cite{lopezcardona2025brain} report converging fMRI--language model evidence consistent with shared representational structure between brains and large models. If this hypothesis is correct in its strong form, two sufficiently capable encoders should embed semantically related items into spaces that are isometric up to a global transformation, and a structural method such as GW should be able to recover the correspondence.

However, recent reassessments complicate this pictures. Gröger et al. \cite{groger2026aristotelian} introduce a permutation-based null calibration and find that the apparent global convergence largely disappears once trivial baselines are accounted for, whereas agreement on \emph{local neighbourhood} structure persists across modalities. They term this refinement the Aristotelian Representation Hypothesis. Ciernik et al. \cite{ciernik2024objective} further show that the consistency of representational similarity across datasets is itself objective-dependent, with self-supervised vision objectives generalising more robustly than image-text contrastive ones.

\section{Approach: Gromov-Wasserstein Optimal Transport}

The technical instrument we adopt is the Gromov-Wasserstein (GW) distance \cite{memoli2011gromov,peyre2016gromov}, an extension of optimal transport to settings in which the source and target distributions live in incomparable metric spaces. Unlike the classical Wasserstein distance, which requires a meaningful cost between points across spaces, GW compares only \emph{intra-space} pairwise distances: the alignment that minimises GW preserves, as far as possible, the relational structure of each space without ever computing a cross-modal cost. The framework therefore handles distributions of arbitrary, and potentially mismatched, dimensionality - a CLIP image embedding in $\mathbb{R}^{512}$ may be aligned against a BERT text embedding in $\mathbb{R}^{768}$ without any shared coordinate system or projection.

Gromov-Wasserstein objective have proven effective in improving a variety of cross-modal tasks.
GW alignment of word embedding spaces was introduced by Alvarez-Melis and Jaakkola \cite{alvarez2018gromov}, who showed that purely structural correspondences recover bilingual lexicons competitive with adversarial methods such as MUSE \cite{conneau2018word} and Wasserstein--Procrustes \cite{grave2019unsupervised}. \hl{The same machinery has since been deployed in single-cell multi-omics integration (ensure this section is right)} \cite{demetci2022gromov}, in graph matching with simultaneously learned node embeddings \cite{xu2019gromov}, in joint embedding of heterogeneous metric spaces under unbalanced transport \cite{beier2025joint}, and in cross-brain and cross-model alignment in computational neuroscience, where the GWTune toolbox of Sasaki et al.\ \cite{sasaki2025gwtune} reveals fine-grained one-to-one, group-to-group, and shifted correspondences that supervised similarity analyses overlook. These successes establish GW as a credible candidate for unsupervised alignment of heterogeneous, dimensionality-mismatched, independently produced distributions - exactly the description of cross-dataset audio and image embeddings - but they leave open the question of whether it can recover \emph{semantic} correspondences, not merely structural ones, between large-scale embedding spaces produced by independent foundation models.

\section{The Structure and Alignment Tradeoff}

Pure GW captures geometry but not semantics. We document this gap empirically: GW alignment between independently trained image and text encoders achieves strong distance correlation (Pearson $r = 0.604$) yet weak retrieval ($R@1 = 0.026$ versus $0.323$ for CLIP). The geometric structure that GW recovers is genuine, but it is also ambiguous under permutation: many couplings minimise the GW objective up to a relabelling of points that respects the local distance graph but scrambles their semantic identity. Two clusters of cats and dogs may be correctly grouped while their cross-modal alignments are entirely swapped, and a purely structural objective has no way to distinguish the correct matching from the swapped one.

The remainder of this thesis traces a controlled progression from this pure unsupervised baseline towards an operationally useful recipe. We first incorporate feature information through Fused Gromov--Wasserstein \cite{vayer2020fused}, which interpolates between Wasserstein (feature-only) and GW (structure-only) costs through a single hyperparameter $\alpha$. Because no native cross-modal cost exists in our setting, we construct one through a bridge: a small ridge regression \cite{hoerl1970ridge} fitted on a limited set of labelled pairs maps one space into the other and supplies the Wasserstein term. This produces a hybrid in which a handful of ground-truth correspondences anchor the orientation of the alignment while GW continues to enforce structural consistency on the unlabelled bulk of the data, breaking the permutation ambiguity that limits pure GW. Methods that learn linear or near-linear transformations between latent spaces from minimal alignment information \cite{moschella2023relative,maiorca2023latent} suggest that this anchoring does not need to be expensive; we make this precise by quantifying exactly how much better the model gets with every single labeled pair of data you add.

For modality combinations in which no bridge model exists - the cross-dataset image-to-audio setting that motivates this work - we further investigate two routes. The first constructs caption-mediated pseudo-pairs and trains a direct alignment; the second composes two independently estimated same-dataset plans through a shared text bridge. The latter, a transitive recipe of the form $T_{\text{ia}} = \mathrm{rownorm}(T_{\text{it}} \cdot B \cdot T_{\text{ta}}^\top)$, achieves substantially higher held-out retrieval, indicating that geometry can be carried across modality pairs that were never directly observed together. \hl{This text-mediated composition is conceptually adjacent to the C-MCR/Ex-MCR (is this worth mentioning?)} family of contrastive bridges \cite{wang2023cmcr,wang2023exmcr}, but ours is purely transport-based and requires no contrastive training of a third joint model.

\section{Research Questions}

The thesis is organised around five research questions that progress from existence to mechanism to practicality.

\begin{enumerate}
  \item \textbf{Genuine structural correspondence.} Does Gromov--Wasserstein discover real geometric correspondences between independently trained embedding spaces of different modalities, beyond what would be produced by a random matching? We answer this through cluster normalised mutual information, Pearson distance correlation, $k$-nearest-neighbour overlap, and triplet consistency, each benchmarked against a permuted baseline.
  \cite{huh2024platonic,groger2026aristotelian}.
  \item \textbf{Modality generality.} Does GW generalise across modality combinations? We evaluate image--text, audio--text, text--text cross-lingual, audio--audio, image--image, and cross-dataset image--audio.
  \item \textbf{Minimal supervision and permutation ambiguity.} Can a small number of labelled pairs, combined with FGW, break the permutation ambiguity that limits pure GW, and at what supervision budget does this become competitive with training a dedicated cross-modal model?
\end{enumerate}

\section{Contributions}

The thesis makes the following contributions:

\begin{enumerate}
  \item A systematic measurement of GW structural alignment quality on genuinely independent embedding spaces, accompanied by a comprehensive geometric evaluation benchmarked against permuted baselines.
  \item Empirical quantification of the structural--semantic gap. Pure GW achieves strong geometric agreement (Pearson $r = 0.604$) but weak retrieval ($R@1 = 0.026$ versus CLIP's $0.323$), the first systematic documentation of this separation in a cross-modal setting.
  \item An empirical link between optimal transport theory and the Platonic Representation Hypothesis. Scaling encoders from $43$M to $2.15$B parameters yields monotonic improvement in alignment quality, with OpenCLIP-G surpassing the supervised CLIP baseline ($R@10 = 0.190$ versus $0.182$). The result is consistent with the Aristotelian refinement \cite{groger2026aristotelian} that local neighbourhood structure, rather than global geometry, drives convergence.
  \item A recipe for FGW with minimal supervision: $2$--$20\%$ labelled pairs combined with ridge regression \cite{hoerl1970ridge} as a practical alternative to training cross-modal models. A train/test split of $R@10$ separates the rows on which ridge was fitted (correct by construction) from unseen rows, and shows that the structural FGW term acts as a regulariser on the latter when one of the two modalities is text. At $\alpha = 0.7$ we observe a small but consistent lift over the pure-Wasserstein $\alpha = 0$ baseline (image--text mean $+0.022$, audio--text mean $+0.014$, best single configuration $+0.052$).
  \item Two recipes for cross-dataset image$\,\leftrightarrow\,$audio retrieval, where no bridge model exists: (i) a direct caption-bridged recipe trained on greedy caption-cosine pseudo-pairs that memorises its training rows but does not generalise (held-out $R@10 \le 0.08$ across all $\alpha$); and (ii) a transitive composition that chains two same-dataset plans through a softmax text bridge, $T_{\text{ia}} = \mathrm{rownorm}(T_{\text{it}} \cdot B \cdot T_{\text{ta}}^\top)$, reaching $0.453$ held-out $R@10$ at $350$ ground-truth pairs per leg - a $5.7\times$ improvement over the direct recipe at matched supervision. Cluster-level routing saturates at $2$--$10\%$ supervision for same-dataset pairs, $25$--$50\times$ less than what sample-level $R@10$ requires.
\end{enumerate}
